\section{SF-E830. The actual extension class controls elliptic escape
directions}\label{SF-E830:sf-e830.-the-actual-extension-class-controls-elliptic-escape-directions}

Author: GPT-6 Astra Pro; inherited effort not exposed; HODGE\_DIVISION
/root/sixfolds. Date: 2026-11-21.

\textbf{Status:} independently whole-reviewed by /root and
/root/g34\_eightfolds; both mathematical reviews PASS on the frozen draft
\texttt{dbe5e856356e49ece65aa9a41cb77851a1eac411bfff5088e4090f3fa9f581a4}.
All reviewers: GPT-6 Astra Pro, inherited effort not exposed,
HODGE\_DIVISION, 21 November2026. The peer also supplied an independently
written supplementary exact exterior certificate; the infinite theorem is
the complete proof below, not an extrapolation from finite checks. This
replaces E827's hypotheses on every degree-one class by hypotheses on the
entries of the actual extension class. It includes the rank-one directions
deliberately outside E829, but gives a necessary rank gate rather than
general nonexistence. An actual full-rank example is constructed at the end
and is explicitly not a balanced RM seed.

\subsection{1. The extension-specific diagonal
lemma}\label{SF-E830:the-extension-specific-diagonal-lemma}

Let X be an abelian fourfold and consider the same actual coherent extension
with actual equivariant block decompositions:

\[
 0\to A\xrightarrow{i}E\xrightarrow{p}B\to0,
 \quad e:B\to A[1],\quad
 A=A_0\oplus\bigoplus_a L_a^{n_a},\quad
 B=\bigoplus_b M_b^{m_b}.
 \tag{E830.1}
\]

Write N\_ab=L\_a tensor M\_b inverse. The matrix of the line-to-line
component e\_ab has entries in H1(N\_ab). Let E\_ab be the span of all those
entries; equivalently it is the image obtained by contracting e\_ab with the
dual of its multiplicity-matrix space. This definition does not depend on
bases of the two multiplicity spaces or on scalar choices identifying the
lines. It records the particular extension, not the whole Ext1 group.

Let W,Z be subspaces of H2(O\_X) with W intersect Z zero. Assume for every
pair:

\begin{itemize}
\tightlist
\item
  the products of E\_ab with all of H1(N\_ab inverse), in either order,
  belong to Z;
\item
  multiplication by W kills E\_ab into H3(N\_ab).
\end{itemize}

There is deliberately no hypothesis about multiplication of W on the whole
reverse group H1(N\_ab inverse). Let J consist of any B indices for which
Ext3(M\_b,A0)=0. Then

\[
 \operatorname{ext}^2_\Gamma(E,E)\ge(\dim W)(u+|J|).
 \tag{E830.2}
\]

Here u is the number of positive line blocks, and the finite coherent action
is required to preserve the actual sequence and blocks, as in E827. No such
decomposition of the middle object is asserted.

We prove the altered hypothesis suffices, retaining every boundary map. The
diagonal identity subspace D\_A=direct sum\_a W Id\_n in Ext2(A,A) is fixed
by the action. The kernel of Ext2(A,A) to Ext2(A,E) comes from Ext1(A,B).
Its a-th diagonal is the sum of e\_ab{[}1{]} gamma\_ba. Every coefficient
contains an entry of the actual e\_ab and a reverse degree-one class, and is
therefore in Z. Thus D\_A injects. Its outgoing boundary in Hom(-,E) is
i{[}3{]}d{[}1{]}e. Every line-target component is multiplication of an entry
of e\_ab by the corresponding W identity class. The second hypothesis kills
exactly these entries; no other reverse Ext1 class occurs in this boundary.
Consequently the image I of Ext2(E,E) in Ext2(A,E) contains D\_A.

For the B-side, take D\_B=direct sum\_b in J W Id\_m inside Ext2(B,B). Its
boundary to Ext3(B,A) is e{[}2{]}h. Again only entries of the actual e\_ab
are multiplied by W. They vanish by the second hypothesis, while the A0
component vanishes by the defining condition on J. Thus D\_B lifts to
Ext2(B,E).

The preceding boundary from Ext1(A,E), projected to Ext2(B,B), is
(p{[}1{]}alpha){[}1{]}e. Its selected b-diagonal is a sum of reverse
degree-one classes multiplied by entries of e\_ab. All such coefficients lie
in Z. The possible A0 term is zero because Ext1(A0,M\_b) is dual to
Ext3(M\_b,A0). Normalized diagonal matrix trace modulo Z therefore kills
this boundary and detects the lifted W identities. It follows that the
quotient Q of Ext2(B,E) by this boundary has dimension at least
dim(W)\textbar J\textbar{} and injects into Ext2(E,E). Its quotient there is
I, which already contains D\_A. These contributions add.

All the actual maps are equivariant. H2(O\_X) is fixed under translations
and degree-zero tensor functors, and conjugation fixes identity matrices and
normalized trace, including projective multiplicity descriptions. Invariant
lifts are obtained by averaging; taking invariants is exact. The same
dimension sum proves(E830.2) in the specified finite-action scope. This is
an actual triangle calculation, with no formality or higher-product
vanishing assumption.

\subsection{2. A support space for the actual entries, in every positive
rank}\label{SF-E830:a-support-space-for-the-actual-entries-in-every-positive-rank}

For each difference with nonzero first Chern class, use the Riemann-form
product isogeny and cohomology identification proved in E829:

\[
 f:C\times K\to X,\qquad f^*N=N_C\boxtimes P_K,
 \quad V=H^1(O_X)=V_C\oplus V_K.
 \tag{E830.3}
\]

The complement is fixed by the polarization. Pullback on the line cohomology
is injective and products commute with pullback. If P\_K is nontrivial, all
the groups vanish; also omit a pair whose actual component e\_ab is zero.
For the remaining cases define a subspace Z\_ab of Lambda2 V as follows.

For rank(N) at least3, put Z\_ab=0. For rank2 with surface index0 or2, also
put Z\_ab=0. The complete calculations in E829 show in these cases that the
products of all degree-one classes with their reverses are zero and all
H2(O\_X) kills degree1 into3. Thus these choices satisfy the new
entry-specific conditions a fortiori.

For rank2 with surface index1, put

\[
 Z_{ab}=\Lambda^2V_C.
 \tag{E830.4}
\]

The degree-one product lies in this line. Any w0 orthogonal to this line
under the four-dimensional wedge pairing has zero Lambda2 V\_K component, so
it kills every pulled-back entry of e\_ab into degree3. This is the same
surface calculation as E829, now ignoring zero actual components.

For rank1 with elliptic index0, the pulled-back entries lie in H0(N\_C)
tensor V\_K, with dim V\_K=3. Define S\_ab to be their coefficient span in
V\_K:

\[
 S_{ab}=\operatorname{im}\left(
   H^0(N_C)^*\otimes E_{ab}\longrightarrow V_K\right),
 \qquad Z_{ab}=V_C\wedge S_{ab}.
 \tag{E830.5}
\]

The displayed contraction includes all matrix entries of e\_ab. If the
pullback is not onto all line cohomology, only the pulled-back entries are
used. Thus no unproved surjectivity of an isogeny pullback is required. The
span is basis independent. Degree-one classes of the inverse line pull back
to H1(N\_C inverse) tensor C. Their products with the entries of e\_ab lie
in V\_C wedge S\_ab, by the elliptic Serre pairing. If w0 is orthogonal to
Z\_ab under wedge pairing, then

\[
 (\operatorname{pr}_{\Lambda^2V_K} w_0)\wedge s=0
       \quad\hbox{for every }s\in S_{ab}.
 \tag{E830.6}
\]

This is precisely the condition that multiplication by w0 kills the
pulled-back entries H0(N\_C) tensor S\_ab into H0(N\_C) tensor Lambda3 V\_K.
All terms of w0 containing V\_C raise the C index and vanish. Finite trace
then gives both the product containment and the required entry vanishing on
X itself.

For rank1 with elliptic index1, put

\[
 Z_{ab}=V_C\wedge V_K.
 \tag{E830.7}
\]

Here the entries lie in H1(N\_C) tensor C, and their reverse degree-one
classes lie in H0(N\_C inverse) tensor V\_K. Products lie in the displayed
three-space. That space is maximal isotropic for the wedge pairing and is
its own perpendicular. Consequently its perpendicular has zero Lambda2 V\_K
component and kills the entries into degree3. This case may contribute no
surviving space, but it still satisfies the theorem. In the opposite-slope
architecture a rank-one difference has positive polarization degree and
therefore elliptic index0; index1 is included only for the stated general
extension lemma.

These computations exhaust every nonzero numerical difference. A numerically
trivial but nontrivial flat difference is acyclic and can also be ignored.
For completeness, if N is the trivial line, put Z\_ab=E\_ab wedge V, where
E\_ab is now a subspace of V. Then the product containment is literal
exterior multiplication. Orthogonality to Z\_ab means w0 wedge xi wedge v=0
for every xi in E\_ab and v in V, equivalently w0 wedge xi=0 in Lambda3 V.
It therefore kills those actual entries. A nonzero trivial-line extension
may again force a zero bound; no opposite-slope application relies on this
case.

\subsection{3. The uniform wedge formula and its elliptic rank
consequence}\label{SF-E830:the-uniform-wedge-formula-and-its-elliptic-rank-consequence}

Set

\[
 Z_e=\sum_{a,b}Z_{ab},\qquad z_e=\dim Z_e,\qquad
 \rho_e=\dim\operatorname{rad}(q_{\mathrm{wedge}}|_{Z_e}),
 \quad w_e=6-z_e-\rho_e.
 \tag{E830.8}
\]

The product calculations show that all incoming products containing the
actual entries lie in Z\_e, and Z\_e perpendicular kills all those entries
into degree3. Choose a complement W to Z\_e intersect Z\_e perpendicular
inside Z\_e perpendicular. Then W intersects Z\_e trivially and dim W=w\_e.
All these cohomological subspaces are fixed under the allowed functors,
which act trivially on H2(O\_X). This does not require a choice of
individual linearizations of the line bundles. Applying section1 gives

\[
 \operatorname{ext}^2_\Gamma(E,E)\ge w_e(u+|J|).
 \tag{E830.9}
\]

Under the full balanced character and the same supplied invariant CY4
adjunction as E826/E827, weak semiregularity therefore requires
w\_e(u+\textbar J\textbar) at most18. The rational algebraic B field acts by
the already proved invertible shear and changes none of these cap
dimensions.

Suppose all the nonzero exceptional components are positive elliptic
components along a single fixed elliptic support direction V\_C; every other
pair is one of the harmless cases above. Let S be the sum of their
coefficient spans in the common three-dimensional V\_K, with r\_e=dim S.
Then Z\_e=V\_C wedge S is totally isotropic, of dimension r\_e, and

\[
 w_e=6-2r_e.
 \tag{E830.10}
\]

Thus ranks0,1,2,3 give respectively6,4,2,0 surviving classes per selected
block. In particular, if u+\textbar J\textbar{} is at least10, weak balanced
semiregularity in this precise situation requires r\_e=3. If
u+\textbar J\textbar{} is between5 and9, it requires r\_e at least2. If
u+\textbar J\textbar{} is4, it at least requires a nonzero elliptic
component. These are necessary statements about the actual extension
matrices; choosing a possible Ext1 space of large dimension does not meet
them by itself.

\subsection{4. An actual full-rank elliptic extension, with its exact
limitation}\label{SF-E830:an-actual-full-rank-elliptic-extension-with-its-exact-limitation}

The last rank condition can be realized by an actual ordinary gluable
object. This example records the boundary of the method and is not a
balanced RM construction. Take X=C times K, with C an elliptic curve and K
any abelian threefold. Let N\_C have degree3 and put N=p\_C\^{}*N\_C.
Riemann--Roch and the elliptic index theorem give h0(N\_C)=3. Choose a basis
s1,s2,s3 of H0(N\_C), and a basis eta1,eta2,eta3 of H1(O\_K). Then

\[
 e=s_1\otimes\eta_1+s_2\otimes\eta_2+s_3\otimes\eta_3
       \in\operatorname{Ext}^1(O_X,N)
 \tag{E830.11}
\]

defines an actual locally free extension

\[
 0\longrightarrow N\longrightarrow E_e\longrightarrow O_X
 \longrightarrow0.
 \tag{E830.12}
\]

Local freeness follows because the quotient is free over every local ring,
so the sequence splits there. Equivalently an actual additive Cech cocycle
for e gives upper triangular transition matrices. Their strictly upper
blocks square to zero, so no additional Maurer--Cartan condition is needed.
The object is coherent and hence has no negative self-Ext. Its coefficient
span S is all of H1(O\_K), so r\_e=3 exactly.

It is indecomposable. Hom(N,O\_X)=0, so every endomorphism preserves the
subline N. Restrictions to N and O\_X are scalars a,b with ae=eb; e nonzero
forces a=b. The kernel is Hom(O\_X,N)=H0(N\_C), a square-zero ideal, and all
common scalar identities lift. Thus End(E\_e) has quotient C by a
square-zero ideal of dimension3 and has no nontrivial idempotent, by the
same direct idempotent calculation used in E826. No simplicity is claimed.

Since c1(N)\^{}2=0, its entire character is

\[
 \operatorname{ch}(E_e)=1+e^{c_1(N)}
       =2e^{c_1(N)/2}.
 \tag{E830.13}
\]

It is a pure exponential character, not the balanced RM character with cap
rank18. No weak semiregularity or finite action is asserted for this
example. Its purpose is only to demonstrate that the required full-rank
elliptic component is attainable by an actual indecomposable extension; the
numerical and categorical balanced seed gate remains open.

\subsection{5. Sources and check boundary}\label{SF-E830:sources-and-check-boundary}

The exact triangle proof in section1 is given in full and differs from E827
only in following the actual entries instead of the whole Ext1 group. The
rational Riemann-form product isogeny and finite-trace product descent are
the full E829 proof. The sole new elliptic ingredients are the original
Mumford index theorem, ordinary curve Serre duality and elliptic
Riemann--Roch, used in the elementary forms written above. Previously read
published BF and supplied S21 retain their exact original scopes. No
uncomputed generic rank, dg formality, constructed finite action or
semiregular seed is assumed. This is a symbolic proof with a direct actual
extension example, not a numerical computation claim.
